\section{Addition on \texorpdfstring{$\mathbb{N}$}{N}}

\begin{topicbox}{Recursive Definition of Addition}
\begin{exposition}
Addition on a Peano system is defined from the successor primitive by
recursion. For each fixed left input $m$, the map $n \mapsto m+n$ is generated
from the initial value $S(m)$ and the successor operator on the same Peano
system. This definition makes addition a derived operation rather than an
independent axiom.
\end{exposition}

\begin{definitionbox}{Definition (Addition on a Peano System)}
\begin{definition}[Addition on a Peano System]
\label{def:addition-on-peano-system}
Let $(P,S,1)$ be a Peano system. For each $m \in P$, define the function
\[
f_m : P \to P
\]
to be the unique function satisfying
\[
f_m(1)=S(m)
\qquad\text{and}\qquad
\forall n \in P\;(f_m(S(n))=S(f_m(n))).
\]
For $m,n \in P$, define
\[
m+n \coloneqq f_m(n).
\]
\end{definition}
\end{definitionbox}

\begin{remark*}[Standard quantified statement]
\[
\begin{aligned}
&\text{Let } (P,S,1) \text{ be a Peano system. For each } m \in P,
\text{ let } f_m : P \to P \text{ satisfy}\\
&f_m(1)=S(m)
\qquad\text{and}\qquad
\forall n \in P\;(f_m(S(n))=S(f_m(n))).\\
&\text{Then for all } m,n \in P,\quad m+n \coloneqq f_m(n).
\end{aligned}
\]
\end{remark*}

\begin{remark*}[Negated quantified statement]
\[
\begin{aligned}
&\text{Let } (P,S,1) \text{ be a Peano system. For each } m \in P,
\text{ let } f_m : P \to P \text{ satisfy}\\
&f_m(1)=S(m)
\qquad\text{and}\qquad
\forall n \in P\;(f_m(S(n))=S(f_m(n))).\\
&\text{Then for all } m,n,x \in P,\quad x \neq m+n
\Longleftrightarrow
x \neq f_m(n).
\end{aligned}
\]
\end{remark*}

\begin{remark*}[Failure modes]
The definition fails only when the designated value does not agree with the
unique recursively generated function attached to the fixed left input $m$.
Once the recursive map $f_m$ is fixed, every incorrect candidate for $m+n$ is
simply a value of $P$ different from $f_m(n)$.
\end{remark*}

\begin{remark*}[Failure mode decomposition]
\[
\begin{aligned}
&\text{Let } (P,S,1) \text{ be a Peano system. For each } m \in P,
\text{ let } f_m : P \to P \text{ satisfy}\\
&f_m(1)=S(m)
\qquad\text{and}\qquad
\forall n \in P\;(f_m(S(n))=S(f_m(n))).\\
&x \neq m+n
\Longleftrightarrow
\underbrace{x \neq f_m(n)}_{\text{candidate value disagrees with the recursive output}}.
\end{aligned}
\]
\end{remark*}

\begin{remark*}[Interpretation]
Addition is introduced by fixing the left input and recursively generating the
right-input map. The value $m+1$ is set to be the successor of $m$, and each
further successor step in the right input advances the output by one more
successor. This makes addition repeated succession.
\end{remark*}

\begin{dependencies}
\begin{itemize}
  \item \hyperref[def:peano-system]{Peano System}
  \item \hyperref[thm:peano-iterator-theorem]{Peano Iterator Theorem}
\end{itemize}
\end{dependencies}

\begin{theorembox}{Theorem (Addition Is Well-Defined on a Peano System)}
\begin{theorem}[Addition Is Well-Defined on a Peano System]
\label{thm:addition-well-defined-on-peano-system}
Let $(P,S,1)$ be a Peano system. For every $m \in P$, there exists a unique
function
\[
f_m : P \to P
\]
such that
\[
f_m(1)=S(m)
\qquad\text{and}\qquad
\forall n \in P\;(f_m(S(n))=S(f_m(n))).
\]
Consequently the rule
\[
(m,n) \longmapsto m+n \coloneqq f_m(n)
\]
defines a unique binary operation on $P$.
\hyperref[prf:addition-well-defined-on-peano-system]{Go to proof.}
\end{theorem}
\end{theorembox}

\begin{remark*}[Standard quantified statement]
\[
\begin{aligned}
&\text{Let } (P,S,1) \text{ be a Peano system.}\\
&\forall m \in P\;\exists! f_m : P \to P\;
\Bigl(f_m(1)=S(m) \land
\forall n \in P\;(f_m(S(n))=S(f_m(n)))\Bigr).
\end{aligned}
\]
\end{remark*}

\begin{remark*}[Interpretation]
The theorem shows that the recursive clauses for addition determine exactly one
right-input map for each fixed left input. The addition operation is therefore
not assumed; it is constructed from the successor map by recursion.
\end{remark*}

\begin{remark*}[Source alignment]
This is the house-style well-definedness form of addition obtained by applying
the Peano recursion theorem in the manner of Feferman's construction of
arithmetic operations \citep{FefermanNumberSystems1964}.
\end{remark*}

\begin{dependencies}
\begin{itemize}
  \item \hyperref[def:addition-on-peano-system]{Addition on a Peano System}
  \item \hyperref[thm:peano-iterator-theorem]{Peano Iterator Theorem}
\end{itemize}
\end{dependencies}
\end{topicbox}

\begin{topicbox}{Recursive Clauses as Reusable Equalities}
\begin{exposition}
The recursive definition of addition is given by the base value and successor
step for the right-input map $n \mapsto m+n$. Those clauses are part of the
definition, but they become substantially more useful once isolated as named
theorems. Later algebraic arguments repeatedly cite these equalities rather
than reopening the recursive construction each time.
\end{exposition}

\begin{theorembox}{Theorem (Addition With \texorpdfstring{$1$}{1})}
\begin{theorem}[Addition With $1$]
\label{thm:addition-with-one}
Let $(P,S,1)$ be a Peano system. For every $m \in P$,
\[
m+1=S(m).
\]
\hyperref[prf:addition-with-one]{Go to proof.}
\end{theorem}
\end{theorembox}

\begin{remark*}[Standard quantified statement]
\[
\begin{aligned}
&\text{Let } (P,S,1) \text{ be a Peano system.}\\
&\forall m \in P\;(m+1=S(m)).
\end{aligned}
\]
\end{remark*}

\begin{remark*}[Interpretation]
This theorem records the base clause of the recursive definition in the
operation notation itself. Adding the distinguished first element applies one
successor step to the left input.
\end{remark*}

\begin{dependencies}
\begin{itemize}
  \item \hyperref[def:addition-on-peano-system]{Addition on a Peano System}
\end{itemize}
\end{dependencies}

\begin{theorembox}{Theorem (Addition Successor on the Right)}
\begin{theorem}[Addition Successor on the Right]
\label{thm:addition-successor-on-right}
Let $(P,S,1)$ be a Peano system. For every $m,n \in P$,
\[
m+S(n)=S(m+n).
\]
\hyperref[prf:addition-successor-on-right]{Go to proof.}
\end{theorem}
\end{theorembox}

\begin{remark*}[Standard quantified statement]
\[
\begin{aligned}
&\text{Let } (P,S,1) \text{ be a Peano system.}\\
&\forall m,n \in P\;(m+S(n)=S(m+n)).
\end{aligned}
\]
\end{remark*}

\begin{remark*}[Interpretation]
This theorem records the recursive step of addition in operation notation.
Moving one successor step forward in the right input moves the sum one
successor step forward as well.
\end{remark*}

\begin{dependencies}
\begin{itemize}
  \item \hyperref[def:addition-on-peano-system]{Addition on a Peano System}
\end{itemize}
\end{dependencies}
\end{topicbox}

\begin{topicbox}{Left Input Initialization}
\begin{exposition}
The recursive definition controls addition when the right input is advanced by
successor. To use addition symmetrically, the first structural task is to
understand what happens when the distinguished first element appears on the
left. The next theorem shows that left addition by $1$ also produces a single
successor step.
\end{exposition}

\begin{theorembox}{Theorem (One Plus \texorpdfstring{$n$}{n})}
\begin{theorem}[One Plus $n$]
\label{thm:one-plus-n}
Let $(P,S,1)$ be a Peano system. For every $n \in P$,
\[
1+n=S(n).
\]
\hyperref[prf:one-plus-n]{Go to proof.}
\end{theorem}
\end{theorembox}

\begin{remark*}[Standard quantified statement]
\[
\begin{aligned}
&\text{Let } (P,S,1) \text{ be a Peano system.}\\
&\forall n \in P\;(1+n=S(n)).
\end{aligned}
\]
\end{remark*}

\begin{remark*}[Interpretation]
The theorem shows that the distinguished first element acts on either side as
a single successor step. This is the first sign that the recursively defined
operation is not tied permanently to the asymmetry of its definition.
\end{remark*}

\begin{dependencies}
\begin{itemize}
  \item \hyperref[thm:addition-with-one]{Addition With $1$}
  \item \hyperref[thm:addition-successor-on-right]{Addition Successor on the Right}
  \item \hyperref[thm:induction-principle-for-peano-system]{Induction Principle for a Peano System}
\end{itemize}
\end{dependencies}
\end{topicbox}

\begin{topicbox}{Associativity and Commutativity}
\begin{exposition}
Once the recursive clauses are available as reusable theorems, addition can be
shown to satisfy the two fundamental algebraic symmetries of a binary
operation. Associativity makes repeated addition unambiguous, and
commutativity shows that the recursively defined operation does not depend on
which input is treated as the evolving one.
\end{exposition}

\begin{theorembox}{Theorem (Addition Is Associative)}
\begin{theorem}[Addition Is Associative]
\label{thm:addition-associative}
Let $(P,S,1)$ be a Peano system. For every $a,b,c \in P$,
\[
(a+b)+c=a+(b+c).
\]
\hyperref[prf:addition-associative]{Go to proof.}
\end{theorem}
\end{theorembox}

\begin{remark*}[Standard quantified statement]
\[
\begin{aligned}
&\text{Let } (P,S,1) \text{ be a Peano system.}\\
&\forall a,b,c \in P\;((a+b)+c=a+(b+c)).
\end{aligned}
\]
\end{remark*}

\begin{remark*}[Interpretation]
Associativity shows that successive additions can be regrouped without
changing the value. This makes longer sums structurally coherent and prepares
addition for later interaction with order and multiplication.
\end{remark*}

\begin{dependencies}
\begin{itemize}
  \item \hyperref[thm:addition-with-one]{Addition With $1$}
  \item \hyperref[thm:addition-successor-on-right]{Addition Successor on the Right}
  \item \hyperref[thm:induction-principle-for-peano-system]{Induction Principle for a Peano System}
\end{itemize}
\end{dependencies}

\begin{theorembox}{Theorem (Addition Is Commutative)}
\begin{theorem}[Addition Is Commutative]
\label{thm:addition-commutative}
Let $(P,S,1)$ be a Peano system. For every $a,b \in P$,
\[
a+b=b+a.
\]
\hyperref[prf:addition-commutative]{Go to proof.}
\end{theorem}
\end{theorembox}

\begin{remark*}[Standard quantified statement]
\[
\begin{aligned}
&\text{Let } (P,S,1) \text{ be a Peano system.}\\
&\forall a,b \in P\;(a+b=b+a).
\end{aligned}
\]
\end{remark*}

\begin{remark*}[Interpretation]
Although addition is defined by recursion on the right input, the resulting
operation is symmetric in its two arguments. Commutativity shows that the
apparent asymmetry of the defining clauses is an artifact of construction
rather than a feature of the operation itself.
\end{remark*}

\begin{dependencies}
\begin{itemize}
  \item \hyperref[thm:addition-with-one]{Addition With $1$}
  \item \hyperref[thm:addition-successor-on-right]{Addition Successor on the Right}
  \item \hyperref[thm:one-plus-n]{One Plus $n$}
  \item \hyperref[thm:addition-associative]{Addition Is Associative}
  \item \hyperref[thm:induction-principle-for-peano-system]{Induction Principle for a Peano System}
\end{itemize}
\end{dependencies}
\end{topicbox}

\begin{topicbox}{Cancellation Laws}
\begin{exposition}
The next structural property shows that addition preserves enough information
to allow one to remove a common summand from an equation. These cancellation
laws are the first substitute for subtraction inside a Peano system and are
indispensable when order is later defined through additive comparison.
\end{exposition}

\begin{theorembox}{Theorem (Left Cancellation for Addition)}
\begin{theorem}[Left Cancellation for Addition]
\label{thm:addition-left-cancellation}
Let $(P,S,1)$ be a Peano system. For every $a,b,c \in P$,
\[
a+b=a+c
\Longrightarrow
b=c.
\]
\hyperref[prf:addition-left-cancellation]{Go to proof.}
\end{theorem}
\end{theorembox}

\begin{remark*}[Standard quantified statement]
\[
\begin{aligned}
&\text{Let } (P,S,1) \text{ be a Peano system.}\\
&\forall a,b,c \in P\;\bigl((a+b=a+c) \Longrightarrow b=c\bigr).
\end{aligned}
\]
\end{remark*}

\begin{remark*}[Contrapositive quantified statement]
\[
\begin{aligned}
&\text{Let } (P,S,1) \text{ be a Peano system.}\\
&\forall a,b,c \in P\;\bigl(b \neq c \Longrightarrow a+b \neq a+c\bigr).
\end{aligned}
\]
\end{remark*}

\begin{remark*}[Interpretation]
Left cancellation shows that adding the same element to two inputs cannot
erase a difference between them. This is the algebraic core that later makes
additive definitions of order behave well.
\end{remark*}

\begin{dependencies}
\begin{itemize}
  \item \hyperref[ax:peano-successor-injective]{Injectivity of Successor}
  \item \hyperref[thm:addition-with-one]{Addition With $1$}
  \item \hyperref[thm:addition-successor-on-right]{Addition Successor on the Right}
  \item \hyperref[thm:one-plus-n]{One Plus $n$}
  \item \hyperref[thm:addition-associative]{Addition Is Associative}
  \item \hyperref[thm:addition-commutative]{Addition Is Commutative}
  \item \hyperref[thm:induction-principle-for-peano-system]{Induction Principle for a Peano System}
\end{itemize}
\end{dependencies}

\begin{theorembox}{Theorem (Right Cancellation for Addition)}
\begin{theorem}[Right Cancellation for Addition]
\label{thm:addition-right-cancellation}
Let $(P,S,1)$ be a Peano system. For every $a,b,c \in P$,
\[
b+a=c+a
\Longrightarrow
b=c.
\]
\hyperref[prf:addition-right-cancellation]{Go to proof.}
\end{theorem}
\end{theorembox}

\begin{remark*}[Standard quantified statement]
\[
\begin{aligned}
&\text{Let } (P,S,1) \text{ be a Peano system.}\\
&\forall a,b,c \in P\;\bigl((b+a=c+a) \Longrightarrow b=c\bigr).
\end{aligned}
\]
\end{remark*}

\begin{remark*}[Contrapositive quantified statement]
\[
\begin{aligned}
&\text{Let } (P,S,1) \text{ be a Peano system.}\\
&\forall a,b,c \in P\;\bigl(b \neq c \Longrightarrow b+a \neq c+a\bigr).
\end{aligned}
\]
\end{remark*}

\begin{remark*}[Interpretation]
Right cancellation is the companion to left cancellation. Once commutativity
is known, the placement of the common summand no longer matters, so equality
can be tested after adding the same element on either side.
\end{remark*}

\begin{dependencies}
\begin{itemize}
  \item \hyperref[thm:addition-commutative]{Addition Is Commutative}
  \item \hyperref[thm:addition-left-cancellation]{Left Cancellation for Addition}
\end{itemize}
\end{dependencies}
\end{topicbox}

\begin{topicbox}{Additive Non-Return}
\begin{exposition}
Addition on $\mathbb{N}$ always moves forward by at least one successor step
in its right input. The next theorem records the resulting non-return law:
adding a natural number to $m$ never returns the added natural number itself.
\end{exposition}

\begin{theorembox}{Theorem (No Natural Number Equals Itself Plus a Natural Number)}
\begin{theorem}[No Natural Number Equals Itself Plus a Natural Number]
\label{thm:no-natural-number-equals-itself-plus-a-natural-number}
Let $(P,S,1)$ be a Peano system. For every $m,n \in P$,
\[
n \neq m+n.
\]
\hyperref[prf:no-natural-number-equals-itself-plus-a-natural-number]{Go to proof.}
\end{theorem}
\end{theorembox}

\begin{remark*}[Standard quantified statement]
\[
\begin{aligned}
&\text{Let } (P,S,1) \text{ be a Peano system.}\\
&\forall m,n \in P\;(n \neq m+n).
\end{aligned}
\]
\end{remark*}

\begin{remark*}[Interpretation]
Because every element of $P$ is at least one successor step, the sum $m+n$
lies strictly beyond $n$. The theorem rules out additive cycles and is a
standard tool for order arguments on the natural numbers.
\end{remark*}

\begin{remark*}[Source alignment]
This is the house-style form of the additive non-return theorem in
\citep{FefermanNumberSystems1964}.
\end{remark*}

\begin{dependencies}
\begin{itemize}
  \item \hyperref[thm:addition-with-one]{Addition With $1$}
  \item \hyperref[thm:addition-successor-on-right]{Addition Successor on the Right}
  \item \hyperref[thm:no-object-is-its-own-successor]{No Object Is Its Own Successor}
  \item \hyperref[thm:induction-principle-for-peano-system]{Induction Principle for a Peano System}
\end{itemize}
\end{dependencies}
\end{topicbox}
